\documentclass[10pt,a4paper]{amsart}
\usepackage[margin=19mm]{geometry}
\usepackage{amsmath,amssymb,mathtools}
\usepackage[scr=rsfs,cal=euler]{mathalfa}
\usepackage{xfrac}
\usepackage[unicode,hidelinks,bookmarks=false]{hyperref}
\newcommand{\ie}{i.e.\ }
\newcommand{\cf}{cf.\ }
\newcommand{\resp}{respectively\ }
\newcommand{\Subsection}{\subsection}
\providecommand{\nobreakdash}{\nobreak\hbox{-}}
\setlength{\emergencystretch}{2em}
\multlinegap0pt
\usepackage{amssymb}
\usepackage{xcolor}
\usepackage[normalem]{ulem}
\usepackage{tikz}
\usepackage{float}
\newtheorem{proposition}{Proposition}
\title{Weakly $\sqrt{J}U$ Rings}
\author{Zari Vesali Mahmood, Ahmad Moussavi, Peter Danchev}
\date{Source: arXiv:2602.14610v1 (CC BY 4.0)}
\begin{document}
\maketitle

\noindent\emph{Source and licence.} Weakly $\sqrt{J}U$ Rings. Version arXiv:2602.14610v1 (CC BY 4.0), distributed by arXiv under the Creative Commons Attribution 4.0 International licence, http://creativecommons.org/licenses/by/4.0/, as recorded in the arXiv licence metadata for this exact version and shown on the abstract page https://arxiv.org/abs/2602.14610v1. Full licence notice: \texttt{licenses/arxiv-2602.14610-CC-BY-4.0.txt}.\par\smallskip

\noindent\textit{Editorial scope.} Counted unit: the article's proposition matrix (source0.tex lines 283--285). The article's complete original proof of it is delivered with it (source0.tex lines 287--299, one \texttt{proof} environment of 13 lines). No mathematics is written, repaired or re-derived: the statement and the proof are the article's own lines, in the article's own notation, and no retained line is substituted or re-flowed.  Retained uncounted as the source-local context the proof needs: lines 371--373.  Nothing was omitted from any retained span, so the package carries no omission record.  No retained line contains a citation, so the wrapper carries no bibliography and no bibliography entry.  No retained line contains a figure environment, an image-inclusion command or drawing code, so no image is delivered and the package declares no asset dependency.  Editorial formatting only, in addition: an AMS \texttt{amsart} preamble that restates the article's own declarations and macros used by the retained lines, namely the environments \texttt{proposition} on separate counters, together with the standard packages those lines need.  The retained environments print in this document as Proposition 1 in order of appearance, whereas the article prints the same items with the numbering of its own full document.  The complete native source of the article as distributed in the arXiv e-print for this exact version is bundled unchanged as \texttt{sources/194727/source0.tex}.
\begin{proposition}\label{corner} For any \( e \in {\rm Id}(R) \), if the ring \( R \) is $W\sqrt{J}U$, then
the corner ring \( eRe \) is also $W\sqrt{J}U$.
\end{proposition}

\begin{proposition}\label{matrix} For any non-zero ring \( R \) and any natural number \( n \geq 2 \), the full
matrix ring \( {\rm M}_n(R) \) is not a $W\sqrt{J}U$ ring.
\end{proposition}

\begin{proof} Since \( {\rm M}_2(R) \) is isomorphic to a corner ring of \( {\rm M}_n(R) \), in conjunction with Proposition~\ref{corner} listed and proved below, it suffices to show that \( {\rm M}_2(R) \) is not a $W\sqrt{J}U$ ring.

To that goal, consider the matrix unit $A = \begin{pmatrix} 0 & 1 \\ 1 & 1 \end{pmatrix}$. Now, one observes that
\[ A + I =
\begin{pmatrix} 0 & 1 \\ 1 & 1 \end{pmatrix} + \begin{pmatrix} 1 & 0 \\ 0 & 1 \end{pmatrix} = \begin{pmatrix} 1
& 1 \\ 1 & 2 \end{pmatrix}, \] which is a unit and hence cannot lie in $\sqrt{J(R)}$.

Furthermore, one sees that \[ A - I = \begin{pmatrix} 0 & 1 \\ 1 & 1 \end{pmatrix} - \begin{pmatrix} 1 & 0 \\ 0 & 1
\end{pmatrix} = \begin{pmatrix} -1 & 1 \\ 1 & 0 \end{pmatrix}, \] which is a unit too, and so it is not in
$\sqrt{J(R)}$ either.

Consequently, one concludes that \( {\rm M}_2(R) \) is not a $W\sqrt{J}U$ ring, as pursued.
\end{proof}

\end{document}
